% =============================================================
%  Section: Discussion
% =============================================================

This section interprets the numbers reported in Section~\ref{sec:results} in
the light of the bug investigation of Section~\ref{sec:bug-investigation}
and the architectural choices of Section~\ref{sec:impl:model}, distinguishes
between the failure modes that are inherent to the YOLOv1 formulation and
the ones that are specific to our training regime, lists the lessons that
generalise beyond this project, and sketches the modifications that would
most directly raise the absolute performance within the same family of
detectors.

% -------------------------------------------------------------
\subsection{Why the absolute mAP is low}
\label{sec:disc:why-low}

The headline figure of $\mathrm{mAP}@0.5 = 0.0415$ on the COCO test split is
roughly a factor of fifteen below the $63.4$ that
\citet{redmon2016you} report on VOC~2007. Four factors, each independent of
implementation correctness, contribute roughly multiplicatively to this gap.

\paragraph{COCO is a substantially harder benchmark than VOC.} The class
count rises from $20$ to $80$, which by itself dilutes the average across
more failure-prone categories. More importantly, the object size
distribution shifts: COCO contains a long tail of small objects (a
substantial fraction of annotations are smaller than the $64\!\times\!64$
footprint of a single cell in our $S=7$ grid, see
Section~\ref{sec:results:perclass}), and many of its scenes are
foreground-clustered in ways that VOC is not (multiple cutlery instances on
dining tables, multiple people on streets, multiple sports objects in
foreground photographs). The single-object-per-cell constraint of YOLOv1 is
much more frequently activated on COCO than on VOC.

\paragraph{Training is roughly half the duration of the paper.} V6 trains for
$60$ epochs total; \citet{redmon2016you} train for $135$. The cosine schedule
in our run reaches its floor by epoch $\sim\!45$, so additional epochs at
that learning rate are unlikely to extract much more from the current
configuration; a fresh schedule at lower peak rate could.

\paragraph{The backbone is smaller than the paper's.} ResNet18 carries
$\sim\!11.2$~M parameters in the convolutional trunk;
Darknet24~\citep{redmon2016you} carries $\sim\!39$~M. ResNet18 also has a
shallower receptive field at the relevant spatial resolution, which limits
the contextual information available to a single-cell prediction. We chose
ResNet18 for tractability (a pretrained checkpoint ships with torchvision and
fits a 50-epoch budget on a single H200), not because it is optimal for the
task. The trade-off paid in headline mAP is the obvious one.

\paragraph{Data augmentation is minimal.} \citet{redmon2016you} apply random
scaling, translation, exposure and saturation jitter, in addition to the
standard horizontal flip. Our pipeline applies only resize. Detection
benchmarks are particularly sensitive to scale augmentation because every
training sample teaches the network about a single absolute object size, and
without augmentation the distribution of trained object sizes is exactly the
distribution in the training set. Mosaic augmentation, introduced by later
YOLO variants, would also help by densifying the number of objects per
training image.

\medskip

\noindent The combined effect of these four factors is consistent with an
order-of-magnitude penalty on mAP, which is what we observe. Crucially, none
of the four is an implementation bug; they are deliberate or pragmatic
choices made within the scope of this project. The bugs of
Section~\ref{sec:bug-investigation} are an entirely separate axis, and they
have been verified to be absent in V5--V6 by the round-trip checks discussed
in that section and reinforced by tests added together with the fixes.

% -------------------------------------------------------------
\subsection{Why specific classes fail}
\label{sec:disc:class-failures}

The per-class breakdown of Section~\ref{sec:results:perclass} identified four
recurring failure modes among the twenty COCO classes with $\mathrm{AP}=0$:
sub-cell objects, clustered tableware, rare classes and sports objects in
human-dominated foregrounds. Each of these is an interaction between an
architectural choice and a property of the dataset.

\paragraph{Sub-cell objects and the $S=7$ grid.} Each grid cell covers a
$64\!\times\!64$ region of the $448\!\times\!448$ input. An object whose
bounding box is smaller than this in both dimensions cannot be precisely
localised by the regression head: the centre regression target $(x, y)$
quantises to within the cell, and the $\sqrt{w}, \sqrt{h}$ regression
targets are pulled towards values whose dynamic range the network has
limited opportunity to learn. Toothbrushes, hair driers, scissors and remote
controls fall into this regime in a substantial fraction of their COCO
appearances.

\paragraph{Clustered tableware and the one-object-per-cell rule.} When
multiple objects of the same class share a cell in the ground truth, only
one is retained as the regression target; the rest are silently dropped.
Fork/knife/spoon scenes are the canonical case, but the same mechanism
applies to multiple birds on a wire, multiple bottles on a shelf, or
multiple cars in a parking lot. At evaluation time, the missed instances
count as false negatives, depressing recall and therefore AP.

\paragraph{Rare classes.} Hair drier (5 ground-truth instances in the test
split), toaster (9), parking meter (59) and similar low-frequency classes
have not seen enough training examples to learn a robust template. The
exception is bear (67 instances in test, $\mathrm{AP}=0.17$): rarity alone
is not disqualifying, provided the silhouette is distinctive enough to be
learnt from a few hundred training examples.

\paragraph{Foreground competition.} Sports objects (ball, frisbee, baseball
bat, baseball glove, snowboard) frequently appear in scenes whose dominant
object is a person. The class loss for the responsible cell is then in
direct competition with the much stronger and more consistent gradient from
the co-located person instance. The result is a cell that learns to
classify the scene as ``person'' regardless of which class actually owns the
cell's centre.

\medskip

\noindent None of these failure modes are unique to our implementation;
\citet{redmon2016you} discusses several of them as known limitations of
YOLOv1, particularly the inability to handle clusters of small objects.
What is specific to our setting is the \emph{degree} to which they hurt the
metric: COCO triggers all four failure modes more often than VOC does, and
our smaller backbone leaves less spare capacity to overcome them.

% -------------------------------------------------------------
\subsection{Lessons that generalise beyond this project}
\label{sec:disc:lessons}

The bug investigation in Section~\ref{sec:bug-investigation} exposes a
handful of failure modes whose details are specific to YOLO but whose shape
generalises to any deep-learning project that involves structured outputs.
We summarise them here.

\paragraph{A passing test is not evidence of correctness.} The pre-fix unit
tests for \texttt{YoloLoss}, \texttt{convert\_cellboxes} and the dataset all
passed; the headline result was nonetheless zero. The tests covered tensor
shapes, dtype, value ranges and trivial sanity checks (``the loss runs
without crashing'', ``the output has the expected dimension'') but did not
encode any quantitative invariant of what the function is supposed to
compute. A test of the form ``encode a known box, decode it, assert the
recovered box matches'' is structurally different from a test of the form
``the output has the right shape'', and only the former would have caught
Bug~2. The same gap exists in many production codebases, and is worth
auditing systematically rather than only when a result is mysterious.

\paragraph{Coordinate round-trips deserve explicit checks.} Detection
pipelines move bounding boxes through three different coordinate systems
(absolute pixels, image-normalised, cell-relative for $x,y$) and one
quantisation step (assignment to a grid cell). Any one of these steps can
silently distort the output; the cumulative effect is hard to debug from
the model's predictions alone. An explicit round-trip---encode a known box
through the dataset code, decode it through the post-processing code, and
assert equality---would have detected Bug~2 in a single line. We added such
a test together with the fix.

\paragraph{In-place tensor ops can sever autograd.} The in-place
\texttt{box\_predictions[..., 2:4] = sign(.) * sqrt(.)} pattern shipped in
an earlier version of the loss was syntactically valid, did not raise an
exception, returned the correct loss value, and silently broke the backward
pass through the $w, h$ pathway. PyTorch generally accepts this pattern;
it only objects when the result tensor is also needed by a later operation
in the same forward pass. The detection is essentially impossible without
either a unit test for gradient flow through the offending pathway, or the
deliberate habit of expressing per-cell transforms with \texttt{torch.cat}
rather than indexed assignment.

\paragraph{A converging loss is not evidence of a learning model.}
The most counter-intuitive observation of the bug investigation is that
V1--V4 produced training and validation losses that decreased monotonically
across epochs, plateaued at sensible-looking values (around $1.0$ for V4),
and gave a strong impression that the model was learning. They were not.
The loss they were minimising was an arithmetically valid quantity that the
network could indeed reduce, but it was not the YOLO loss; it was a
projection of the YOLO loss onto a subset of the output channels that did
not include the box-coordinate slots. A useful generalisation: when a model
that should detect things detects nothing, the question
``is the loss the loss?'' should rank above
``are the hyper-parameters tuned?''.

% -------------------------------------------------------------
\subsection{Future work}
\label{sec:disc:future}

The most direct routes to a higher absolute mAP, ordered by expected impact
within the constraints of the YOLOv1 architectural family, are:

\begin{enumerate}
    \item \emph{Larger backbone with a comparable budget.} Replacing
        ResNet18 with ResNet50 roughly triples the trunk parameter count and
        approximately doubles the convolutional depth. On ImageNet
        classification this is worth several percentage points of top-1
        accuracy; the corresponding gain on COCO detection should be larger
        than that of a single hyper-parameter sweep at the current backbone.
    \item \emph{Finer grid ($S = 14$).} A $14\!\times\!14$ grid at the
        current $448\!\times\!448$ input would shrink the per-cell footprint
        to $32\!\times\!32$ pixels, which is below the threshold at which the
        sub-cell-object failure mode of Section~\ref{sec:disc:class-failures}
        triggers for most COCO classes. This requires either dropping the
        adapter's downsampling stride or doubling the input resolution to
        $896\!\times\!896$ (at quadratic cost in inference time).
    \item \emph{Mosaic augmentation.} The mosaic technique introduced by
        later YOLO variants stitches four random training images into a
        single sample, multiplying the number of objects per training
        example by four and indirectly training the network on object
        instances in a wider range of contexts. This is cheap to implement
        and has been reported to be worth several mAP points on COCO with
        small detectors.
    \item \emph{Anchor boxes.} \citet{redmon2018yolov3} replace the
        ``responsible box'' assignment of YOLOv1 with a small set of
        per-grid-cell anchor priors (clustered $k$-means on the training-set
        boxes) and convert the box regression from direct $(x, y, w, h)$
        prediction to a delta against each anchor. This decouples the
        regression target from the per-cell competition between boxes and
        substantially improves localisation. It would require changes in the
        dataset, the loss and the post-processing simultaneously, and is
        therefore the largest item on this list.
    \item \emph{Longer training, more aggressive LR.} V5--V6 stop at the
        cosine floor after $60$ total epochs. A fresh schedule of
        $150$~epochs with a peak rate of $5\!\times\!10^{-4}$ (intermediate
        between V1's $10^{-3}$ and V4's $10^{-4}$) and the existing
        gradient-clip threshold of~$10$ would likely yield additional
        improvement, although less than any of the four items above.
\end{enumerate}

\noindent Items 1, 3 and 5 are continuous improvements; item 2 is an
architectural change but within the YOLOv1 framework; item 4 effectively
constitutes YOLOv2 and would be a substantially different system.
